%% ma: L12-B05-0033
%% lop: 12
%% chuong: 1
%% bai: 5
%% dang: 12.5.23
%% loai: TLN
%% muc: VD
%% dap_an: "25,2"
%% nguon: "HK1-DeOn_1_1.pdf (M10, Đề ôn tập học kì 1 Toán 12, Đề số 2), Phần III Câu 4"
%% kien_thuc: "Bài 5 (ứng dụng đạo hàm giải quyết vấn đề thực tiễn)"
%% trang_thai: cho-duyet
%% ly_do: "Dạng mới đề xuất 12.5.23"
\begin{ex}
Người ta muốn làm một cái bồn chứa nước dạng hình hộp chữ nhật không có nắp đậy, đáy thùng có chiều dài gấp ba lần chiều rộng và có thể tích $18\,000$ lít. Để giảm chi phí, người ta cần phải thiết kế sao cho tổng diện tích các mặt của bồn chứa nước là nhỏ nhất. Tính chi phí thấp nhất (đơn vị tính triệu đồng) để sản xuất ra một cái bồn. Biết rằng giá vật liệu là $400$ nghìn đồng/$\text{m}^2$ và giá thiết kế, thi công, hoàn thiện cái bồn là $300$ nghìn đồng/$\text{m}^2$.
\shortans{25,2}
\loigiai{Gọi chiều rộng đáy $x$ (m), chiều dài $3x$, chiều cao $h$: $3x^2h=18\Rightarrow h=\dfrac6{x^2}$.

Diện tích $S(x)=3x^2+2(xh+3xh)=3x^2+\dfrac{48}x$, $S'(x)=6x-\dfrac{48}{x^2}=0\Leftrightarrow x=2$; $S_{\min}=S(2)=36$ (m$^2$).

Chi phí thấp nhất: $36\cdot(400+300)=25\,200$ nghìn đồng $=25{,}2$ triệu đồng.}
\end{ex}
